\section{Introduction}\label{sec:introduction}
The Erd\H{o}s--R\'enyi graph $G(n,p)$ is one of the standard models of random graphs, defined as a graph on $n$ vertices in which each edge is present independently with probability $p$.
For a fixed graph $H$, let $X_H$ be the number of injective homomorphisms from
$H$ to $G(n,p)$.  Thus, writing
$(n)_k=n(n-1)\cdots(n-k+1)$,
\[
  \mathbb{E}[X_H]=(n)_{|V(H)|}p^{|E(H)|}.
\]
Our primary interest is the following two questions raised by Chatterjee and Varadhan~\cite{ChatterjeeVaradhanLDP} in the context of the upper-tail large deviation problem: (1) what is the probability that $X_H$ is at least $(1+\delta)\mathbb{E}[X_H]$ for some fixed $\delta>0$?; (2) what is the typical structure of $G(n,p)$ conditioned on this event?


This upper-tail problem has been studied extensively; see, for example,
\cite{BhattacharyaDembo2021,ChatterjeeMissingLog,DeMarcoKahnTriangles, Janson2004uppertail, JansonRucinskiInfamous,Raz2020Cycles, Sileikis2020Stars}.
However, even when $H$ is a triangle, the problem is difficult and not
completely understood.  Existing results take different forms in the sparse
and dense regimes.  For the sparse results below, $\delta>0$ is fixed while
$n\to\infty$ and $p=p(n)\to0$.
The polynomial concentration results of Kim and Vu~\cite{KimVu2000} shows that the probability that $X_{K_3}\ge(1+\delta)\mathbb{E}[X_{K_3}]$ lies between $\exp(-O(n^2p^2 \log(1/p)))$ and $\exp(-\Omega(n^2p^2))$. 
Chatterjee~\cite{ChatterjeeMissingLog} and, independently, DeMarco and
Kahn~\cite{DeMarcoKahnTriangles} determined the correct order when $p\ge \log n/n$ by showing that 
\[
  \mathbb{P}\bigl(X_{K_3}\ge(1+\delta)\mathbb{E}[X_{K_3}]\bigr)
  =\exp\bigl(-\Theta_\delta(n^2p^2\log(1/p))\bigr).
\]
Here and throughout the paper, $\log$ denotes the natural logarithm.  These results
determine the order of the exponent, but not its leading constant.  Chatterjee
and Dembo~\cite{ChatterjeeDemboNonlinear} subsequently proved a nonlinear
large-deviation principle which, when $p$ does not decay too quickly, reduces
the upper-tail problem to a weighted-graph relative-entropy variational
problem.
Lubetzky and Zhao~\cite{LubetzkyZhaoSparse} solved this variational problem for
triangles and, by combining their solution with the Chatterjee--Dembo
principle, obtained, for $n^{-1/42}\log n\le p\ll1$,
\[
  \mathbb{P}\bigl(X_{K_3}\ge(1+\delta)\mathbb{E}[X_{K_3}]\bigr)
  =
  \exp\left[
    -\left(
      \min\left\{\frac{\delta^{2/3}}2,\frac{\delta}{3}\right\}
      +o(1)
    \right)n^2p^2\log(1/p)
  \right].
\]
This determines the exact leading constant.  It applies, however, over a
narrower range of $p$, since $n^{-1/42}\log n\allowbreak\gg(\log n)/n$.
Afterward, a series of works
established sharp upper-tail asymptotics for broader classes of graphs \(H\)
and in progressively sparser regimes; see
\cite{Augeri2020, BasakBasu2023, BhattacharyaGangulyLubetzkyZhao,CookDembo2020,CookDemboPham2024, Eldan2018}.

In the dense regime, primarily considered in this paper, $p$ is a fixed constant.
In this case, it is convenient to formulate the upper-tail event in
terms of homomorphism densities.  For a graph $G$ on $n$ vertices, define
\[
  t(H,G):=\frac{\hom(H,G)}{n^{|V(H)|}},
\]
where $\hom(H,G)$ is the number of graph homomorphisms from $H$ to $G$.
Since $H$ is fixed, the number of non-injective homomorphisms is
$O_H(n^{|V(H)|-1})$, and hence, for $G=G(n,p)$,
\[
  t(H,G(n,p))=\frac{X_H}{n^{|V(H)|}}+O_H(n^{-1}).
\]
Thus, the homomorphism density $t(H,G(n,p))$ differs from the normalized
injective count $X_H/n^{|V(H)|}$ by only $O_H(n^{-1})$.  We therefore
use the following homomorphism-density version of the upper-tail event.  For
$p<r<1$, write
\[
  \mathcal E_{H,r}:=
  \{t(H,G(n,p))\ge r^{|E(H)|}\}.
\]
The count event in the first paragraph corresponds asymptotically to the
choice $r^{|E(H)|}=(1+\delta)p^{|E(H)|}$. In this paper,
a central role is played by the relative-entropy function
\[
  J_p(u) = u \log \frac{u}{p} + (1-u) \log \frac{1-u}{1-p},
\]
with the convention $0\log0=0$.  This is the large-deviation cost per edge of
raising the edge density from $p$ to $u$.  In this notation, Lubetzky and
Zhao~\cite{lubetzky2015large} showed that if the point $(r^2, J_p(r))$ lies on
the convex minorant of $x \mapsto J_p(x^{1/2})$, then
\[
  \mathbb{P}(\mathcal E_{K_3,r})
  =\exp(-J_p(r)n^2/2+o(n^2)),
\]
and they extended this result to general regular graphs $H$.


Conditioned on $\mathcal E_{H,r}$, the typical structure of $G(n,p)$ can take one of two forms: either it looks like an Erd\H{o}s--R\'enyi graph with a higher edge density (called the \emph{replica-symmetric phase}), or its edges are clustered (called the \emph{symmetry-breaking phase}).
Chatterjee and Varadhan~\cite{ChatterjeeVaradhanLDP} showed that, depending on the parameters, both structures can occur.
Lubetzky and Zhao~\cite{lubetzky2015large} determined exactly when the replica-symmetric phase occurs and when the symmetry-breaking phase occurs in terms of the function $J_p(u)$ defined above for every $d$-regular graph $H$.
\begin{theorem}[Lubetzky--Zhao criterion, informal version~\cite{lubetzky2015large}]\label{thm:lz-criterion-informal}
  Let $H$ be a $d$-regular graph with $d \geq 2$, and let $0<p<r<1$. If the point $(r^d, J_p(r))$ lies on the convex minorant of the function $x\mapsto J_p(x^{1/d})$, then the typical structure of $G(n,p)$ conditioned on $\mathcal E_{H,r}$ is like an Erd\H{o}s--R\'enyi graph with edge density $r$, i.e., $(p, r)$ is in the replica-symmetric phase. Otherwise, the typical structure of $G(n,p)$ conditioned on $\mathcal E_{H,r}$ is clustered, i.e., $(p, r)$ is in the symmetry-breaking phase.
\end{theorem}
See \Cref{thm:lz-criterion} for the formal statement.
The result of Lubetzky and Zhao also yields that the probability of $\mathcal{E}_{H, r}$ is $\exp(-J_p(r)n^2/2 + o(n^2))$ when $(p, r)$ is in the replica-symmetric phase.
However, in the symmetry-breaking phase, neither the manner in which the
edges cluster nor, in general, the leading exponential decay rate of
$\mathbb P(\mathcal E_{H,r})$ is known.
Lubetzky and Zhao~\cite[Section~7]{lubetzky2015large} raised the question of
whether the conditioned random graph always has a bipodal structure in this
phase.  Informally, this means that the vertex set of $G(n,p)$ can be partitioned
into two parts such that the edges within each part and between the two parts
behave like Erd\H{o}s--R\'enyi graphs, with possibly different edge densities.  Numerical experiments using the
neural graphon representation of Kim, Baek, Lee, and
Yang~\cite{KimBaekLeeYangGraphons} provide further evidence for an
affirmative answer to this question; see \Cref{fig:neural-bipodal-patterns}.

Our main contribution is a step towards resolving this bipodality question
and has two complementary parts.  Away from the exceptional density
$r_*=(d-1)/d$, we prove bipodality and uniqueness throughout a nontrivial
neighbourhood of the phase boundary on the symmetry-breaking side.  At $r_*$,
where this nondegenerate local theory breaks down, we instead find a family of
balanced bipodal optimizers along an analytic curve approaching the phase boundary.
The Lubetzky--Zhao phase boundary can be written as $p=\pc(r)$,
as we establish in \Cref{thm:scalar-lz-boundary}.
Our expansions are naturally expressed in terms of the log-odds
distance to this boundary,
\[
  \lambda(p, r) = \log \frac{1-p}{p} - \log \frac{1-\pc(r)}{\pc(r)},
\]
which vanishes exactly on the boundary and is positive on the symmetry-breaking side $p<\pc(r)$.
\begin{theorem}[Informal main results]\label{thm:main-results-informal}
  Let $H$ be a $d$-regular graph with $d \geq 2$, and let $\pc(r)$ be the Lubetzky--Zhao boundary curve.
  Set $r_*=(d-1)/d$.
  \begin{enumerate}[label=(\alph*)]
    \item Let $r_0\in(0,1)\setminus\{r_*\}$.  There exists an open
    neighbourhood $U$ of $(\pc(r_0),r_0)$ such that, for every
    $(p,r)\in U$ with $p<\pc(r)$, the typical structure of $G(n,p)$
    conditioned on $\mathcal E_{H,r}$ is bipodal.  After labelling the
    smaller block first, its relative size and the three edge densities are
    uniquely determined and depend analytically on $(p,r)$.
    The relative size of the smaller block tends to zero as $(p,r)$
    approaches the phase boundary.

    Furthermore, for $(p,r)\in U$ with $p<\pc(r)$,
    \[
      \mathbb P(\mathcal E_{H,r})
      =\exp\left[-\left(J_p(r)-\frac{\lambda(p,r)^2}{2A_H(r)}
        +O(\lambda(p,r)^3)\right)\frac{n^2}{2}+o(n^2)\right],
    \]
    with the implied constants uniform on $U$, where
    $A_H:(0,1)\setminus\{r_*\}\to(0,\infty)$ is an analytic function that
    depends only on $r$ and $H$.

    \item There exists an analytic curve $h\mapsto(p_h,r_h)$ of
    symmetry-breaking parameters approaching the exceptional boundary point:
    \[
      (p_h,r_h)\longrightarrow(\pc(r_*),r_*)
      \qquad(h\downarrow0).
    \]
    Along this curve, the typical structure of $G(n,p_h)$ conditioned on
    $\mathcal E_{H,r_h}$ is described by a unique nonconstant bipodal graphon
    $W_h$, up to relabelling.  As $h\downarrow0$, both block proportions
    tend to $1/2$, while all three edge densities tend to $r_*$.
  \end{enumerate}
\end{theorem}
The two parts describe qualitatively different mechanisms by which a
nonconstant bipodal optimizer approaches a constant graphon at the phase
boundary.  Away from $r_*$, one pode shrinks to measure zero while the other
occupies almost all of $[0,1]$ and its within-pode density tends to $r$.  The
densities involving the shrinking pode need not tend to $r$, but the regions
on which they occur have vanishing measure.  The optimizer therefore converges
to $W\equiv r$ in $L^1$, and hence in cut distance, but generally not in
$L^\infty$.  At $r_*$, neither pode disappears: their sizes tend to $1/2$, while
both within-pode densities and the cross density tend to $r_*$.  Thus, the
exceptional optimizer converges in $L^\infty$ by homogenization of two
macroscopic podes, rather than by the disappearance of a small pode.
The precise variational statements underlying parts~(a) and~(b) are
\Cref{thm:nonexceptional-optimizers,thm:endpoint-optimizers}, respectively.

\subsection{Chatterjee--Varadhan variational problem}
We now recast the preceding random-graph questions as a variational problem on
graphons.
A \emph{graphon} is a symmetric measurable function $W:[0,1]^2\to[0,1]$.
We write $\W_0$ for the space of all graphons.
Graphons are limit objects for dense graphs; see~\cite{LovaszGraphLimits}.
For a graphon $W$, define the \emph{edge density}, \emph{homomorphism density}, and \emph{relative entropy} by 
\begin{align*}
  e(W) & =\int_{[0,1]^2}W(x,y)\,dx\,dy,
  \\
  t(H,W) & =\int_{[0,1]^{|V(H)|}}\prod_{ij\in E(H)}W(x_i,x_j)\prod_{i\in V(H)}dx_i,
  \\
  I_p(W) & =\int_{[0,1]^2}J_p(W(x,y))\,dx\,dy.
\end{align*}
For $0<p<r<1$ and a graph $H$, define the Chatterjee--Varadhan variational problem~\cite{ChatterjeeVaradhanLDP} by
\begin{equation}\label{eq:graphon-variational-problem}
    \inf_{W\in\W_0} I_p(W)\quad \text{subject to } t(H,W)\ge r^{|E(H)|}.
\end{equation}
We denote by $\Phi_H(p,r)$ the optimal value of the above variational problem
and by $\W_H(p,r)$ its set of minimizers.  A graphon is \emph{feasible} for
this problem if it satisfies the constraint
$t(H,W)\ge r^{|E(H)|}$; the collection of all such graphons is the feasible
set.  This set is nonempty because the constant graphon $W\equiv r$ is
feasible.  Standard compactness, continuity, and lower-semicontinuity results
for graphons imply that the infimum is attained, so $\W_H(p,r)$ is nonempty~\cite[Theorem~2.7]{lubetzky2015large}.
The cut distance $\delta_\square$ measures the difference between two graphons
after allowing measure-preserving relabellings; its formal definition, together
with the graphon-topological facts used above, is given in
\Cref{sec:graphons}.  A graph $G$ can be viewed as its associated step graphon
$W_G$, in which case $t(H,G)=t(H,W_G)$.
The following theorem makes precise the connection between the variational
problem, the upper-tail probability, and the conditioned random graph.
\begin{theorem}[Chatterjee--Varadhan large-deviation principle~\cite{ChatterjeeVaradhanLDP}]
\label{thm:graphon-large-deviations}
  For every $0<p<r<1$, we have 
  \[
    \lim_{n \to \infty}\frac{1}{\binom{n}{2}}
    \log \mathbb{P}\bigl(t(H,G(n,p))\ge r^{|E(H)|}\bigr)
    =-\Phi_H(p,r).
  \]
  Furthermore, for any $\eps>0$, there exists $C=C(\eps, p, r, H)>0$ such that, for all sufficiently large $n$,
  \[
    \mathbb{P}\left(
      \min_{W \in \W_H(p,r)}\delta_\square(W_{G(n,p)},W)>\eps
      \ \middle|\ t(H,G(n,p))\ge r^{|E(H)|}
    \right)
    \le e^{-Cn^2}.
  \]
\end{theorem}
The first conclusion gives the leading exponential decay rate of the
upper-tail probability, while the second says that, conditional on the
upper-tail event, the graphon of a typical $G(n,p)$ is close in cut distance
to some optimizer in $\W_H(p,r)$, with probability tending exponentially fast
to one.

In variational terms, the Lubetzky--Zhao criterion stated informally in
\Cref{thm:lz-criterion-informal} takes the following form.
\begin{theorem}[Lubetzky--Zhao criterion~\cite{lubetzky2015large}]\label{thm:lz-criterion}
    Let $H$ be a $d$-regular graph with $d\ge2$ and let $0<p<r<1$.
    If the point $(r^d, J_p(r))$ lies on the convex minorant of the function $x\mapsto J_p(x^{1/d})$, then the unique minimizer of~\eqref{eq:graphon-variational-problem} is the constant graphon $W\equiv r$. Otherwise, the constant graphon is not contained in the set of minimizers $\W_H(p, r)$.
\end{theorem}
A graphon $W$ is called \emph{bipodal} if there is a measurable set $A\subseteq[0,1]$ such that, writing $B=[0,1]\setminus A$, the graphon $W$ is almost everywhere equal to a constant on each of the three sets $A\times A$, $(A\times B)\cup(B\times A)$, and $B\times B$; see \Cref{sec:graphons}.
Thus, in the graphon formulation, the question of Lubetzky and Zhao asks
whether every minimizer of~\eqref{eq:graphon-variational-problem} in the
symmetry-breaking phase is bipodal.  The precise version of our result away
from the exceptional density is as follows.

\begin{theorem}\label{thm:nonexceptional-optimizers}
    Let $H$ be a $d$-regular graph with $d\ge2$.
    If $r_0 \in (0, 1)\setminus\{(d-1)/d\}$, then there exists an open neighbourhood $U$ of $(\pc(r_0), r_0)$ such that for any $(p, r)\in U$ with $p<\pc(r)$, the minimizer of~\eqref{eq:graphon-variational-problem} is unique up to relabelling and bipodal.
    Furthermore, the unique optimizer $W_{p, r}$ satisfies
      \[
    e(W_{p,r})=r-\frac{\lambda(p,r)}{A_H(r)}+O(\lambda(p,r)^2),
    \qquad
    \Phi_H(p,r)=J_p(r)-\frac{\lambda(p,r)^2}{2A_H(r)}
      +O(\lambda(p,r)^3),
  \]
    where $A_H:(0,1) \setminus\{(d-1)/d\} \to (0, \infty)$ is an analytic function.
\end{theorem}
Together with the Chatterjee--Varadhan theorem above, this gives
the asymptotic formula for the upper-tail probability and the
bipodal description of the conditioned random graph stated in
\Cref{thm:main-results-informal}(a).

To explain why the exceptional density requires a separate argument,
consider $H=K_3$, for which $d=2$ and the exceptional density is $1/2$.
Fix $r\ne1/2$, and let $c(p,r)$ denote the size of the smaller
block.  As $p\uparrow\pc(r)$ from the symmetry-breaking side,
\Cref{rmk:bipodal-parameter-expansions} gives the expansion
\[
  c(p,r)
  =\frac{r}{(2r-1)^2A_{K_3}(r)}\lambda(p,r)
   +O_r\bigl(\lambda(p,r)^2\bigr).
\]
Here the $O_r(\lambda(p,r)^2)$ term is a remainder whose absolute value is
at most $C_r\lambda(p,r)^2$ for $p$ sufficiently close to $\pc(r)$.
For this fixed $r$, the coefficient of $\lambda(p,r)$ is finite and
$\lambda(p,r)\to0$, so the expansion implies $c(p,r)\to0$: the smaller
block disappears.

The difficulty arises when $r$ also tends to $1/2$ as the phase boundary is
approached.  The factor $(2r-1)^2$ in the denominator then tends to zero,
and our estimates for $A_{K_3}(r)$ do not determine whether the full
coefficient of $\lambda(p,r)$ remains bounded.
The error constant $C_r$ and the required closeness of $p$ to $\pc(r)$ may
also depend on $r$.  Thus, knowing that $\lambda(p,r)\to0$ is no longer
enough to conclude from this expansion that $c(p,r)\to0$.
The preceding analysis therefore does not settle whether a block still
disappears along approaches to the exceptional boundary point.

The following result provides that separate treatment through the construction of
asymptotically balanced rank-one graphons.  Here, a graphon $W$ is called
\emph{rank-one} if $W(x,y)=f(x)f(y)$ almost everywhere for some measurable
$f:[0,1]\to[0,1]$.  The result supports an affirmative answer to the
bipodality question by producing symmetry-breaking points arbitrarily close to the boundary
whose optimizers are still bipodal, but with different limiting behavior.
\begin{theorem}\label{thm:endpoint-optimizers}
Let $H$ be a $d$-regular graph with $d\ge2$.  There exist $h_0>0$, an
analytic curve $h\mapsto(p_h,r_h)$ for $0<h<h_0$, and nonconstant rank-one
bipodal graphons $W_h$ such that
\[
  (p_h,r_h)\longrightarrow
  \left(\pc\left(\frac{d-1}{d}\right),\frac{d-1}{d}\right)
  \qquad(h\downarrow0).
\]
The graphon $W_h$ satisfies
\[
  t(H,W_h)=r_h^{|E(H)|},
\]
and is the unique minimizer of \eqref{eq:graphon-variational-problem}, up to
measure-preserving relabelling, at $(p_h,r_h)$.  As $h\downarrow0$, the sizes
of its two blocks both converge to $1/2$, while its two within-block densities
and its cross density all converge to $(d-1)/d$.  Consequently, $W_h$
converges in $L^\infty$ to the constant graphon $W\equiv(d-1)/d$.
\end{theorem}

\subsection{Proof strategy}\label{sec:proof-strategy}
The proofs of \Cref{thm:nonexceptional-optimizers,thm:endpoint-optimizers}
use the supporting-line geometry underlying the Lubetzky--Zhao boundary,
but they identify the graphon optimizers by different methods.
Let $r_*=(d-1)/d$ be the exceptional density.

\begin{enumerate}[label=(\roman*),leftmargin=2.2em]
\item \emph{Nonexceptional case (\Cref{thm:nonexceptional-optimizers}).}
For $r\ne r_*$, we first prove a quadratic lower bound for the gap between
the scalar function $x\mapsto J_{\pc(r)}(x^{1/d})$ and its supporting line
at $x=r^d$.  This line also touches the curve at $x=\sm(r)^d$, with
$\sm(r)\ne r$.  In \Cref{sec:lz-boundary}, we show that the gap is at least a
positive constant times the squared distance from $x$ to the nearer of
these two contact points.  The same positive constant works for every $r$
in a fixed closed interval contained in $(0,1)$ that does not contain
$r_*$.  We exclude $r_*$ because there the two contacts merge and the gap
vanishes to fourth order, so a quadratic lower bound no longer holds.

We next reduce the choice of an optimizing graphon to the choice of its
edge density.  In \Cref{sec:local-reduction}, we show that every optimizer
has exactly the required $H$-density $r^{|E(H)|}$, and that its edge density
approaches $r$ from below as $p\uparrow\pc(r)$.  We can therefore write its
edge density as $r-\delta$, where $\delta>0$ is small.  For each such
$\delta$, minimizing $I_p$ among graphons with these two prescribed
densities is equivalent to maximizing graphon entropy.\footnote{The graphon
entropy is $s(W)=-\frac12\int_{[0,1]^2}[W\log W+(1-W)\log(1-W)]\,dx\,dy$,
with the convention $0\log0=0$.}
The theorem of Kenyon, Radin, Ren, and Sadun
\mbox{(KRR--S)}~\cite{kenyon2014entropy} shows that maximizing graphon
entropy under the constraints $e(W)=r-\delta$ and
$t(H,W)=r^{|E(H)|}$ yields a unique bipodal graphon, up to relabelling.
Thus, each $\delta$ selects one bipodal candidate for the original relative
entropy minimization problem.  Write $W_\delta$ for this bipodal candidate.
It is specified by four parameters: the size of its smaller block, the two
within-block edge densities, and the cross-block edge density.  It remains
to choose the $\delta$ for which $W_\delta$ has the least relative entropy.

To make this choice, we compare the candidates with the constant graphon
$W\equiv r$, which has the same $H$-density.  At $p=\pc(r)$, the constant
graphon is optimal, so reducing the edge density by $\delta$ incurs an
extra relative entropy cost.  The scalar gap bound above, combined with
generalized H\"older's inequality, shows that this extra cost is at least
a positive multiple of $\delta^2$.  We also construct a bipodal graphon
with the prescribed densities whose extra cost is at most another
multiple of $\delta^2$.  Since the four parameters of $W_\delta$ depend
analytically on $\delta$ and $r$, these two bounds show
that its extra cost has leading term $\frac12A_H(r)\delta^2$, with a
positive analytic coefficient $A_H(r)$; see \Cref{sec:nondegeneracy}.

To use this boundary calculation when $p<\pc(r)$, we make the dependence
of the excess cost on $p$ explicit.  For fixed $r$ and $\delta$, the
candidate $W_\delta$ does not depend on $p$, since it maximizes graphon
entropy under constraints that involve only $r$ and $\delta$.  The relative
entropy identity \eqref{eq:relative-entropy-identity} gives the exact relation
\[
  I_p(W_\delta)-J_p(r)
  =\bigl[I_{\pc(r)}(W_\delta)-J_{\pc(r)}(r)\bigr]
   -\lambda(p,r)\delta.
\]
The bracketed term is the excess cost at the boundary, analysed in the
previous paragraph; the term $-\lambda(p,r)\delta$ is the correction
caused by changing $p$.  Inserting the boundary expansion therefore gives
the total excess cost at $p<\pc(r)$:
\[
  I_p(W_\delta)-J_p(r)
  =\frac12A_H(r)\delta^2-\lambda(p,r)\delta+O(\delta^3)
  \qquad(\delta\downarrow0,\ r\text{ fixed}).
\]
The linear term favours a reduction in edge density, while the quadratic
term penalizes it.  Balancing the two gives a minimizing value of $\delta$
with leading term $\lambda(p,r)/A_H(r)$.  Strict convexity makes this
minimizing value unique, and hence also makes the graphon optimizer unique
up to relabelling.  The analytic implicit function theorem shows that the
minimizing $\delta$ depends analytically on $(p,r)$.  Evaluating the four
parameters of $W_\delta$ at the minimizing value of $\delta$ determines
the optimizing graphon.  Evaluating $I_p(W_\delta)$ at the same value
gives the expansion of the minimum relative entropy; see
\Cref{sec:local-optimizer-structure}.

\item \emph{Exceptional case (\Cref{thm:endpoint-optimizers}).}
For $r\ne r_*$, the supporting line to the curve
$x\mapsto J_{\pc(r)}(x^{1/d})$ touches it at two distinct points,
$x=r^d$ and $x=\sm(r)^d$.  As $r\to r_*$, these contact points merge at
$x=r_*^d$.  At $r=r_*$, the vertical gap between the curve and its
supporting line vanishes to fourth order at this point.
Since the KRR--S theorem does not cover this density, we develop a separate
argument in \Cref{sec:singular-endpoint}.

We first seek a family of candidate graphons of the product form $W(x,y)=f(x)f(y)$
for the original relative entropy minimization problem.
Such graphons are called \emph{rank-one}.
Here $f$ assigns a weight to each vertex, and the edge probability is the
product of the two vertex weights.  This form simplifies the $H$-density:
since every vertex of $H$ has degree $d$,
\[
  t(H,W)=\left(\int_0^1 f(x)^d\,dx\right)^{|V(H)|}.
\]
Thus, the $H$-density constraint reduces to a condition on a single
integral of $f$.  We use the product form to construct candidate graphons;
we then prove their optimality by comparison with arbitrary graphons
satisfying the required $H$-density constraint.

To determine the candidate graphons, we impose necessary conditions for a
constrained minimum: the derivative of the relative entropy must vanish
under differentiable perturbations that keep the $H$-density fixed.
These are the \emph{stationarity conditions}.  For nonconstant graphons of
this product form with $f$ bounded away from $0$ and $1$, the resulting equations
force $f$ to take exactly two values.  The sets on which $f$ takes these
values become the two blocks of the graphon.  We also impose stationarity
under changes of block size.  We use $h>0$ to denote half the difference
between the two values of $f_h$.  The stationarity conditions for the
graphon values and for changes in block size together determine the two
values of $f_h$, the block sizes, and $p_h$ as analytic functions of $h$
near $0$.
We obtain $W_h(x,y)=f_h(x)f_h(y)$ and define $r_h$ by
$t(H,W_h)=r_h^{|E(H)|}$; see \Cref{sec:rank-one-stationary-family}.
As $h\downarrow0$, the block sizes tend to $1/2$ and the parameter pair
$(p_h,r_h)$ approaches $(\pc(r_*),r_*)$.
We then show that $W_h$ has lower relative entropy, by a quantity of order
$h^4$, than the constant graphon $W\equiv r_h$ with the same $H$-density;
see \Cref{sec:constant-graphon-comparison}.

It remains to show that $W_h$ is the unique optimizer among all graphons.
Take any graphon $W$ satisfying the density constraint
$t(H,W)\ge r_h^{|E(H)|}$ and having relative entropy
$I_{p_h}(W)\le I_{p_h}(W_h)$.  We must prove that $W$ agrees with $W_h$ up to
relabelling.  The fourth-order supporting gap described above provides
the first step: together with the density constraint and this relative entropy
comparison, it gives
\begin{equation}\label{eq:intro-quartic-localization}
  \int_{[0,1]^2}|W(x,y)-r_*|^4\,dx\,dy\le C_d h^4.
\end{equation}
Thus, every such competitor is within $O(h)$ of the constant graphon
$r_*$ in the $L^4$-norm.  The localization estimate
\eqref{eq:intro-quartic-localization}, proved in
\Cref{sec:localization-rank-one}, allows us to carry out the relative
entropy comparison below.

We can now compare such a $W$ with $W_h$ more precisely.  We write
$W(x,y)=f(x)f(y)+E(x,y)$, separating a product part from a remainder $E$.
We choose $f$ so that replacing $W$ by its product part changes the
$H$-density by at most a cubic quantity in the $L^2$-norm of $E$:
\begin{equation}\label{eq:intro-density-error}
  \left|t(H,W)-\left(\int_0^1 f(x)^d\,dx\right)^{|V(H)|}\right|
  \le C_H\|E\|_2^3.
\end{equation}
For this choice of $f$, the localization bound
\eqref{eq:intro-quartic-localization} gives
$\|f-\sqrt{r_*}\|_4=O(h)$ and $\|E\|_2=O(h)$.
These two bounds provide the smallness needed to control the higher-order
terms in the relative entropy comparison.

We then establish a lower bound for the relative entropy difference
$I_{p_h}(W)-I_{p_h}(W_h)$ that accounts both for the remainder $E$
and for deviations of $f$ from the two values
$s_h<t_h$ taken by $f_h$.  Combining this bound with the density
constraint and the cubic estimate
\eqref{eq:intro-density-error} shows that, for sufficiently
small $h$, $I_{p_h}(W)\ge I_{p_h}(W_h)$, with equality only if
\[
  E=0\ \text{a.e.},\qquad
  f(x)\in\{s_h,t_h\}\ \text{a.e.},\qquad
  \int_0^1 f(x)^d\,dx=\int_0^1 f_h(x)^d\,dx.
\]
Our assumption $I_{p_h}(W)\le I_{p_h}(W_h)$ therefore forces all three
conditions.  The first two say that $W$ itself has the product form,
with the same two factor values as $W_h$.  The last condition fixes the
block sizes: since $s_h$ and $t_h$ are distinct, their proportions are
uniquely determined by the integral of $f^d$.
Thus, $W$ agrees with $W_h$ up to relabelling.  This proves optimality
and uniqueness along the constructed curve; see
\Cref{sec:auxiliary-lagrangian,sec:graphon-comparison,sec:endpoint-optimality-proof}.
\end{enumerate}

\paragraph{Organization.}
\Cref{sec:preliminaries} records the graphon preliminaries and the
entropy-maximization results.  The nonexceptional argument occupies
\Cref{sec:lz-boundary,sec:local-reduction,sec:nondegeneracy,sec:local-optimizer-structure},
and \Cref{sec:singular-endpoint} proves the exceptional result.
\Cref{sec:lean-formalization} describes our Lean 4 formalization,
including its scope and assumptions.
\Cref{sec:concluding-remarks} discusses further directions.
Appendices~\ref{app:krrs-analytic-extension}, \ref{app:lz-scalar-geometry},
\ref{app:bipodal-parameter-expansions}, and \ref{app:local-reduction} contain,
respectively, the derivation of the two-sided KRR--S parametrization, the
scalar convex-analysis proofs, the bipodal parameter expansions, and
proofs supporting the one-dimensional reduction.